% ============================================================
\section{Introduction}
\label{sec:intro}
% ============================================================

Deploying object detectors at the edge requires quantization
\cite{jacob2018quant,gholami2022survey}, and
mixed-precision quantization (MPQ)---spending high precision only where
it matters---promises a better accuracy--cost trade-off than uniform
bit-widths. The dominant recipe is two-staged: (i) score the
``sensitivity'' of each parameter group (layer, block, or channel)
independently, once, at full precision or at a fixed reference state;
(ii) solve a budgeted allocation problem over these per-group scores
\cite{hawq,haq,bitsplit,flbq,lampq,blobq,regptq}. A silent assumption
underlies stage (i): that a group's contribution is an
\emph{intrinsic, context-free} property.

This paper shows that in object detection this assumption fails, and we
build an allocation method around its failure.

\smallskip
\noindent\textbf{(1) The problem: precision utility is context-dependent.}
  Using a controlled instrument---fake quantization of the 12 transformer
  blocks of a Swin-T backbone inside Mask R-CNN \cite{maskrcnn,swin}, with
a fully GT-free utility defined by agreement with the full-precision (FP)
teacher detector on 64 unlabeled images---we measure precision utility in
\emph{context}: the marginal value of repairing a block given that other
blocks are already repaired. Three findings emerge: pairwise interactions
average $0.79\times$ the size of singleton effects; conditional and
singleton rankings reverse on $22.6$--$32.0\%$ of within-context pairs
across five calibration seeds (Fig.~\ref{fig:reversal}); and at an
identical repair budget, decisions made from singleton rankings lose
$3.56$ AP with respect to conditional measurement (singleton top-4:
$41/495$; conditional: $11/495$, by exact enumeration).

\smallskip
\noindent\textbf{(2) The method: GT-free conditional allocation.}
We present a \emph{conditional} allocation method: it estimates marginal
utility in the current allocation context with the same instrument,
stabilizes noisy 64-image estimates with an $\varepsilon$-tolerant
split-half criterion over single upgrades and budget-neutral pairwise
exchanges, and selects among several starting allocations by full-64
utility. It requires no labels, no retraining, and only
$\sim\!10^{2}$--$10^{3}$ calibration forwards.

\smallskip
\noindent\textbf{(3) The results: best-or-tied at every budget, two backbones.}
At matched parameter-weighted budgets on COCO 5k \cite{coco}, the method
attains the best-or-tied deployable accuracy at every tested budget
(e.g., $+1.8$ over the strongest reimplemented prior at $4.5$-bit), with
the margins holding when the strongest baseline is granted the same
extended bit-level set. It transfers to a second backbone of a different
size (Swin-S, $+12.5$ at equal budget), and it comes with analyses that
explain \emph{why} it works: a sample-complexity account of why $64$
calibration images suffice, and a reversal-floor law that predicts the
measured reversal rate and yields a design criterion for when singleton
scoring is structurally unreliable.
